\subsection{Communication as action}
\label{sec:communication-as-action}
\slidesrc{Communication: 2--8}

\begin{nbdef}{Communication (slide 2)}
Communication, in general, is the \textbf{intentional} exchange of information brought about
by the production and perception of \textbf{signs} drawn from a \textbf{shared system of
conventional signs}.
\end{nbdef}

Two views of language:
\begin{itemize}[nosep]
  \item \textbf{``classical'' view} (before 1953): language consists of sentences that are
        true or false, as in logic (Part~II);
  \item \textbf{``modern'' view} (after 1953): language is a \textbf{form of action}.
\end{itemize}

\begin{center}
\small
\renewcommand{\arraystretch}{1.2}
\begin{tabularx}{\linewidth}{@{}l l X@{}}
\toprule
\textbf{Who} & \textbf{Work} & \textbf{Key idea} \\
\midrule
Wittgenstein & \emph{Philosophical Investigations} (1953) & The structure of reality
  determines the structure of language. \textbf{Meaning is determined by use}, in the
  context of a ``language game'': words are defined not by the objects they designate, nor
  by the mental representations associated with them, but by how they are used in
  effective, ordinary communication. \\
Austin & \emph{How to Do Things with Words} (1962) & \textbf{In saying something, we do
  something}: a minister saying ``I now pronounce you husband and wife'' does not describe a
  marriage, he \emph{makes} it. \\
Searle & \emph{Speech Act Theory} (1969) & Each \textbf{illocution} can be described by what
  it attempts to do: an assertion is an undertaking that something is the case; a question
  is an attempt to elicit information; thanking is an expression of gratitude. What counts
  is the assumed intent of the speaker. \\
\bottomrule
\end{tabularx}
\captionof{table}{From ``language describes'' to ``language acts'' (slides 3--5).}
\end{center}

\begin{nbintuition}[Locution, illocution, perlocution]
\notinslides
Austin distinguishes three levels in ``Close the window'': the \emph{locution} (the words
uttered), the \emph{illocution} (the act performed \emph{in} saying them: a request), and the
\emph{perlocution} (the effect obtained \emph{by} saying them: the hearer closes the window,
or refuses). Agent communication languages standardise the \emph{illocutionary} level: a
message says explicitly which act it performs (\code{request}, \code{inform}, \ldots).
\end{nbintuition}

\begin{nbfact}{Speech acts achieve the speaker's goals (slide 6)}
\begin{center}
\small
\begin{tabular}{@{}ll@{}}
\toprule
\textbf{Speech act} & \textbf{Example} \\
\midrule
inform & ``There's a pit in front of you'' \\
query & ``Can you see the gold?'' \\
command & ``Pick it up'' \\
promise & ``I'll share the gold with you'' \\
acknowledge & ``OK'' \\
\bottomrule
\end{tabular}
\end{center}
(The examples come from the ``wumpus world'' of Russell and Norvig: agents exploring a cave
with pits and gold.) \textbf{Speech act planning} requires knowledge of the situation, of
semantic and syntactic conventions, and of the hearer's goals, knowledge base and
rationality.
\end{nbfact}

Treating speech as action means that an agent can \emph{plan} its utterances like any other
action: ``I want $H$ to believe $P$; informing $H$ that $P$ has this effect if $H$ trusts me;
so I inform''.

\begin{nbfact}{Stages in communication (slide 7)}
\begin{center}
\begin{tikzpicture}[font=\small, >=Stealth,
  s/.style={draw, thick, rounded corners, text width=3.0cm, align=center, minimum height=0.9cm, fill=#1}]
  \node[s=nbblue!10] (a) at (0,0) {\textbf{Intention}\\$S$ wants to inform $H$ that $P$};
  \node[s=nbblue!10] (b) at (3.6,0) {\textbf{Generation}\\$S$ selects words $W$ to express $P$ in context $C$};
  \node[s=nbblue!10] (c) at (7.2,0) {\textbf{Synthesis}\\$S$ utters words $W$};
  \node[s=nbgreen!10] (d) at (10.8,0) {\textbf{Perception}\\$H$ perceives $W'$ in context $C'$};
  \node[s=nbgreen!10] (e) at (10.8,-2.0) {\textbf{Analysis}\\$H$ infers possible meanings $P_1,\ldots,P_n$};
  \node[s=nbgreen!10] (f) at (7.2,-2.0) {\textbf{Disambiguation}\\$H$ infers intended meaning $P_i$};
  \node[s=nbgreen!10] (g) at (3.6,-2.0) {\textbf{Incorporation}\\$H$ incorporates $P_i$ into its KB};
  \draw[->, thick] (a) -- (b); \draw[->, thick] (b) -- (c); \draw[->, thick, nbred] (c) -- (d);
  \draw[->, thick] (d) -- (e); \draw[->, thick] (e) -- (f); \draw[->, thick] (f) -- (g);
\end{tikzpicture}
\end{center}
Blue: speaker $S$; green: hearer $H$. Note $W'$ and $C'$: what the hearer perceives and its
context may differ from what the speaker uttered and meant.
\end{nbfact}

\begin{nbwarning}[How could this go wrong? (slide 8)]
\begin{itemize}[nosep]
  \item \textbf{Insincerity}: $S$ does not believe $P$ (a lie, or a malicious agent);
  \item \textbf{recognition errors}: the slide's joke ``speech wreck ignition failure'' is
        how a speech recogniser might hear ``speech recognition failure'' ($W' \neq W$);
  \item \textbf{ambiguous utterance}: several meanings $P_1, \ldots, P_n$ remain;
  \item \textbf{differing understanding of the current context} ($C \neq C'$).
\end{itemize}
Software agents avoid most of these problems by \emph{design}: messages are structured, the
act is explicit, and the meaning of the words is fixed by a shared ontology (next
sub-sections). Insincerity remains a real issue in open systems (trust and reputation).
\end{nbwarning}
